\documentclass[../../main.tex]{subfiles}
\begin{document}
\section{全导数与偏导数}\label{section:appC-total-and-partial-derivatives}

对有限维向量空间（的开子集）之间的映射，最一般的导数概念是全导数。

设 $V,W$ 是有限维向量空间，我们可假定它们赋予范数。如果 $U\subset V$ 是开子集，$a\in U$，映射 $F:U\to W$ 如果存在线性映射 $L:V\to W$ 使
\begin{align}\label{align:appC-differentiable}
\lim_{v\to0}\frac{|F(a+v)-F(a)-Lv|}{|v|}=0,
\end{align}
就称它在 $a$ 处\emph{可微}。此式分子中的范数是 $W$ 的范数，分母中的范数是 $V$ 的范数。由于有限维向量空间上的所有范数都等价（习题 B.49），此定义与两种范数的选取无关。

\begin{exercise}\label{exercise:appC-C-1}
假设 $F:U\to W$ 在 $a\in U$ 处可微。证明满足\eqref{align:appC-differentiable} 的线性映射 $L$ 唯一。
\end{exercise}

如果 $F$ 在 $a$ 处可微，满足\eqref{align:appC-differentiable} 的线性映射 $L$ 记作 $\mathrm{D}F(a)$，称为 $F$ 在 $a$ 处的\emph{全导数}。条件\eqref{align:appC-differentiable} 也可写成
\begin{align}\label{align:appC-remainder}
F(a+v)=F(a)+\mathrm{D}F(a)v+R(v),
\end{align}
其中余项 $R(v)=F(a+v)-F(a)-\mathrm{D}F(a)v$ 满足当 $v\to0$ 时 $|R(v)|/|v|\to0$。因而全导数表示 $a$ 附近对 $F(a+v)-F(a)$ 的“最佳线性逼近”。

\begin{exercise}\label{exercise:appC-C-2}
假设 $V,W,X$ 是有限维向量空间，$U\subset V$ 是开子集，$a\in U$，$F,G:U\to W$ 与 $f,g:U\to\mathbb R$ 是映射。证明下列命题。
\begin{enumerate}[(a)]
\item 如果 $F$ 在 $a$ 处可微，则它在 $a$ 处连续。
\item 如果 $F$ 是常值映射，则 $F$ 在 $a$ 处可微且 $\mathrm{D}F(a)=0$。
\item 如果 $F,G$ 在 $a$ 处可微，则 $F+G$ 也可微，且
\begin{align*}
\mathrm{D}(F+G)(a)=\mathrm{D}F(a)+\mathrm{D}G(a).
\end{align*}
\item 如果 $f,g$ 在 $a$ 处可微，则 $fg$ 也可微，且
\begin{align*}
\mathrm{D}(fg)(a)=f(a)\mathrm{D}g(a)+g(a)\mathrm{D}f(a).
\end{align*}
\item 如果 $f,g$ 在 $a$ 处可微且 $g(a)\neq0$，则 $f/g$ 在 $a$ 处可微，且
\begin{align*}
\mathrm{D}(f/g)(a)=\frac{g(a)\mathrm{D}f(a)-f(a)\mathrm{D}g(a)}{g(a)^2}.
\end{align*}
\item 如果 $T:V\to W$ 是线性映射，则 $T$ 在每一点 $v\in V$ 都可微，全导数等于 $T$ 本身：$\mathrm{D}T(v)=T$。
\item 如果 $B:V\times W\to X$ 是双线性映射，则 $B$ 在每一点 $(v,w)\in V\times W$ 都可微，且
\begin{align*}
\mathrm{D}B(v,w)(x,y)=B(v,y)+B(x,w).
\end{align*}
\end{enumerate}
\end{exercise}

\begin{proposition}[全导数的链式法则]\label{proposition:appC-chain-rule}
假设 $V,W,X$ 是有限维向量空间，$U\subset V$ 与 $\widetilde U\subset W$ 是开子集，$F:U\to\widetilde U$ 与 $G:\widetilde U\to X$ 是映射。如果 $F$ 在 $a\in U$ 处可微、$G$ 在 $F(a)\in\widetilde U$ 处可微，则 $G\circ F$ 在 $a$ 处可微，且
\begin{align*}
\mathrm{D}(G\circ F)(a)=\mathrm{D}G(F(a))\circ\mathrm{D}F(a).
\end{align*}
\end{proposition}

\begin{proof}
令 $A=\mathrm{D}F(a)$、$B=\mathrm{D}G(F(a))$。我们需证
\begin{align}\label{align:appC-chain-proof}
\lim_{v\to0}\frac{|G(F(a+v))-G(F(a))-BAv|}{|v|}=0.
\end{align}
写 $b=F(a)$、$w=F(a+v)-F(a)$。代入后可把\eqref{align:appC-chain-proof} 中的商改写为
\begin{align*}
\frac{|G(b+w)-G(b)-BAv|}{|v|}
=\frac{|G(b+w)-G(b)-Bw+Bw-BAv|}{|v|}
\leqslant\frac{|G(b+w)-G(b)-Bw|}{|v|}
+\frac{|B(w-Av)|}{|v|}.
\end{align*}
由于 $A,B$ 线性，习题 B.52 表明存在常数 $C,C'$，使对所有 $x\in V$ 有 $|Ax|\leqslant C|x|$、对所有 $y\in W$ 有 $|By|\leqslant C'|y|$。$F$ 在 $a$ 处可微意味着：对任意 $\varepsilon>0$，只要 $v$ 落在 $0$ 的足够小邻域内，就可保证
\begin{align*}
|w-Av|=|F(a+v)-F(a)-Av|\leqslant\varepsilon|v|.
\end{align*}
此外，当 $v\to0$ 时，由 $F$ 的连续性，$|w|=|F(a+v)-F(a)|\to0$。因此，$G$ 在 $b$ 处可微意味着：必要时再缩小 $|v|$，还可使
\begin{align*}
|G(b+w)-G(b)-Bw|\leqslant\varepsilon|w|.
\end{align*}

综合这些估计，可见当 $|v|$ 足够小时，上式被
\begin{align*}
\frac{\varepsilon|w|}{|v|}+C'\frac{|w-Av|}{|v|}
=\varepsilon\frac{|w-Av+Av|}{|v|}+C'\frac{|w-Av|}{|v|}
\leqslant\varepsilon\frac{|w-Av|}{|v|}+\varepsilon\frac{|Av|}{|v|}+C'\frac{|w-Av|}{|v|}
\leqslant\varepsilon^2+\varepsilon C+C'\varepsilon
\end{align*}
控制，可使之任意小。
\end{proof}

\subsection*{偏导数}

现在我们专门讨论 Euclid 空间之间的映射。设 $U\subset\mathbb R^n$ 是开集，$f:U\to\mathbb R$ 是实值函数。对任意 $a=(a^1,\dots,a^n)\in U$ 与任意 $j\in\{1,\dots,n\}$，$f$ 在 $a$ 处关于 $x^j$ 的\emph{偏导数}定义为固定其他变量时 $f$ 对 $x^j$ 的普通导数：
\begin{align*}
\frac{\partial f}{\partial x^j}(a)=\lim_{h\to0}\frac{f(a^1,\dots,a^j+h,\dots,a^n)-f(a^1,\dots,a^j,\dots,a^n)}{h}
=\lim_{h\to0}\frac{f(a+he_j)-f(a)}{h},
\end{align*}
若极限存在。

更一般地，对向量值函数 $F:U\to\mathbb R^m$，可把 $F(x)$ 的坐标写成 $F(x)=(F^1(x),\dots,F^m(x))$。这定义了 $m$ 个函数 $F^1,\dots,F^m:U\to\mathbb R$，称为 $F$ 的分量函数。$F$ 的偏导数就定义为其分量函数的偏导数 $\partial F^i/\partial x^j$。偏导数所成矩阵 $(\partial F^i/\partial x^j)$ 称为 $F$ 的\emph{Jacobian 矩阵}，其行列式称为 $F$ 的\emph{Jacobian 行列式}。

如果 $F:U\to\mathbb R^m$ 在 $U$ 内每一点处各偏导数都存在，且如此定义的函数 $\partial F^i/\partial x^j:U\to\mathbb R$ 都连续，则称 $F$ 是 $C^1$ 类的或\emph{连续可微}的。在此情形中，可对函数 $\partial F^i/\partial x^j$ 再求导，得到二阶偏导数
\begin{align*}
\frac{\partial^2 F^i}{\partial x^k\partial x^j}=\frac{\partial}{\partial x^k}\left(\frac{\partial F^i}{\partial x^j}\right),
\end{align*}
若它们存在。继续下去就得到更高阶偏导数：$F$ 的 $k$ 阶偏导数就是 $k-1$ 阶偏导数的（一阶）偏导数，当它们存在时。

一般地，如果 $U\subset\mathbb R^n$ 是开子集、$k\geqslant0$，函数 $F:U\to\mathbb R^m$ 如果 $F$ 的所有阶数不超过 $k$ 的偏导数都存在且是 $U$ 上的连续函数，就称它是 $C^k$ 类的或 $k$ 次连续可微的。（因而 $C^0$ 类函数就是连续函数。）由于导数的存在性与连续性都是局部性质，显然 $F$ 是 $C^k$ 的，当且仅当它在 $U$ 中每一点的某邻域内都有该性质。

对每个 $k\geqslant0$ 都是 $C^k$ 类的函数，称为 $C^\infty$ 类的、\emph{光滑}的或\emph{无限可微}的。如果 $U,V$ 是 Euclid 空间的开子集，映射 $F:U\to V$ 如果光滑、双射且其反函数也光滑，就称为\emph{微分同胚}。

链式法则有一个推论值得注意。

\begin{proposition}\label{proposition:appC-diffeo-dim}
假设 $U\subset\mathbb R^n$ 与 $V\subset\mathbb R^m$ 是开子集，$F:U\to V$ 是微分同胚。则 $m=n$，且对每个 $a\in U$，全导数 $\mathrm{D}F(a)$ 可逆，其逆为 $\mathrm{D}(F^{-1})(F(a))$。
\end{proposition}

\begin{proof}
由于 $F^{-1}\circ F=\operatorname{Id}_U$，链式法则蕴含对每个 $a\in U$，
\begin{align}\label{align:appC-diffeo-proof}
\operatorname{Id}_{\mathbb R^n}=\mathrm{D}(\operatorname{Id}_U)(a)=\mathrm{D}(F^{-1}\circ F)(a)=\mathrm{D}(F^{-1})(F(a))\circ\mathrm{D}F(a).
\end{align}
类似地，$F\circ F^{-1}=\operatorname{Id}_V$ 蕴含 $\mathrm{D}F(a)\circ\mathrm{D}(F^{-1})(F(a))$ 是 $\mathbb R^m$ 上的恒等映射。这蕴含 $\mathrm{D}F(a)$ 可逆且逆为 $\mathrm{D}(F^{-1})(F(a))$，因而 $m=n$。
\end{proof}

有时我们需要考虑定义域为 $\mathbb R^n$ 中非开子集的函数的光滑性。如果 $A\subset\mathbb R^n$ 是任意子集，函数 $F:A\to\mathbb R^m$ 如果在每一点的某开邻域上都有光滑延拓，就称它在 $A$ 上\emph{光滑}；更精确地说，对每个 $x\in A$，存在含 $x$ 的开子集 $U_x\subset\mathbb R^n$ 与光滑函数 $\widetilde F:U_x\to\mathbb R^m$，使 $\widetilde F$ 在 $U_x\cap A$ 上与 $F$ 一致。微分同胚的概念以显然方式推广到任意子集：给定任意子集 $A,B\subset\mathbb R^n$，从 $A$ 到 $B$ 的微分同胚是一个光滑双射 $f:A\to B$，且其逆光滑。

我们特别关心实值函数，即陪域为 $\mathbb R$ 的函数。如果 $U\subset\mathbb R^n$ 是开集，$U$ 上所有 $C^k$ 类实值函数所成集合记作 $C^k(U)$，所有光滑实值函数所成集合记作 $C^\infty(U)$。函数的和、常数倍与积按逐点定义：对 $f,g:U\to\mathbb R$ 与 $c\in\mathbb R$，
\begin{align*}
(f+g)(x)=f(x)+g(x),\qquad (cf)(x)=c(f(x)),\qquad (fg)(x)=f(x)g(x).
\end{align*}

\begin{exercise}\label{exercise:appC-C-5}
设 $U\subset\mathbb R^n$ 是开子集，假设 $f,g\in C^\infty(U)$、$c\in\mathbb R$。
\begin{enumerate}[(a)]
\item 证明 $f+g$、$cf$、$fg$ 都光滑。
\item 证明这些运算使 $C^\infty(U)$ 成为交换环，并成为 $\mathbb R$ 上的交换结合代数。
\item 证明如果 $g$ 在 $U$ 上从不为零，则 $f/g$ 光滑。
\end{enumerate}
\end{exercise}

下列重要结果表明：对大多数有意义的函数，求偏导数的次序无关紧要。证明见 [Apo74, Rud76, Str00]。

\begin{proposition}[混合偏导数相等]\label{proposition:appC-mixed-partials}
如果 $U$ 是 $\mathbb R^n$ 的开子集，$F:U\to\mathbb R^m$ 是 $C^2$ 类函数，则 $F$ 的二阶混合偏导数与求导次序无关：
\begin{align*}
\frac{\partial^2 F^i}{\partial x^j\partial x^k}=\frac{\partial^2 F^i}{\partial x^k\partial x^j}.
\end{align*}
\end{proposition}

\begin{corollary}\label{corollary:appC-mixed-partials}
如果 $F:U\to\mathbb R^m$ 光滑，则 $F$ 的任意阶混合偏导数都与求导次序无关。
\end{corollary}

下面我们研究全导数与偏导数之间的关系。设 $U\subset\mathbb R^n$ 是开集，$F:U\to\mathbb R^m$ 在 $a\in U$ 处可微。作为 $\mathbb R^n$ 与 $\mathbb R^m$ 之间的线性映射，$\mathrm{D}F(a)$ 可等同于一个 $m\times n$ 矩阵。下一个命题把该矩阵确认为 $F$ 的 Jacobian。

\begin{proposition}\label{proposition:appC-jacobian}
设 $U\subset\mathbb R^n$ 是开集，假设 $F:U\to\mathbb R^m$ 在 $a\in U$ 处可微。则 $F$ 在 $a$ 处所有偏导数都存在，且 $\mathrm{D}F(a)$ 是矩阵为 $F$ 在 $a$ 处 Jacobian 的线性映射：
\begin{align*}
\mathrm{D}F(a)=\left(\frac{\partial F^j}{\partial x^i}(a)\right).
\end{align*}
\end{proposition}

\begin{proof}
令 $B=\mathrm{D}F(a)$，对使 $a+v\in U$ 的充分小 $v\in\mathbb R^n$，令 $R(v)=F(a+v)-F(a)-Bv$。$F$ 在 $a$ 处可微意味着：当 $v\to0$ 时，向量值函数 $R(v)/|v|$ 的每个分量都趋于零。$F^j$ 在 $a$ 处的第 $i$ 个偏导数若存在，就是
\begin{align*}
\frac{\partial F^j}{\partial x^i}(a)=\lim_{t\to0}\frac{F^j(a+te_i)-F^j(a)}{t}
=\lim_{t\to0}\frac{B^j_i t+R^j(te_i)}{t}
=B^j_i+\lim_{t\to0}\frac{R^j(te_i)}{t}.
\end{align*}
上面商的范数是 $|R^j(te_i)|/|te_i|$，当 $t\to0$ 时趋于零。由此 $\partial F^j/\partial x^i(a)$ 存在且等于所断言的 $B^j_i$。
\end{proof}

\begin{exercise}\label{exercise:appC-C-9}
假设 $U\subset\mathbb R^n$ 是开集。证明函数 $F:U\to\mathbb R^m$ 在 $a\in U$ 处可微，当且仅当其每个分量函数 $F^1,\dots,F^m$ 都在 $a$ 处可微。证明在此情形中
\begin{align*}
\mathrm{D}F(a)=\begin{pmatrix}\mathrm{D}F^1(a)\\ \vdots\\ \mathrm{D}F^m(a)\end{pmatrix}.
\end{align*}
\end{exercise}

前一习题表明：对开区间 $J\subset\mathbb R$，映射 $\gamma:J\to\mathbb R^m$ 可微，当且仅当其分量函数在单变量微积分意义下可微。

下一个命题给出可微性最重要的充分条件；特别地，它表明初等微积分中所有通常的函数都可微。证明见 [Apo74, Rud76, Str00]。

\begin{proposition}\label{proposition:appC-c1-differentiable}
设 $U\subset\mathbb R^n$ 是开集。如果 $F:U\to\mathbb R^m$ 是 $C^1$ 类的，则它在 $U$ 内每一点处可微。
\end{proposition}

对 Euclid 空间之间的函数，链式法则可用偏导数重述。

\begin{corollary}[偏导数的链式法则]\label{corollary:appC-partial-chain}
设 $U\subset\mathbb R^n$ 与 $\widetilde U\subset\mathbb R^m$ 是开子集，令 $x=(x^1,\dots,x^n)$ 表示 $U$ 上的标准坐标、$y=(y^1,\dots,y^m)$ 表示 $\widetilde U$ 上的标准坐标。
\begin{enumerate}[(a)]
\item $C^1$ 函数 $F:U\to\widetilde U$ 与 $G:\widetilde U\to\mathbb R^p$ 的复合仍是 $C^1$ 类的，其偏导数为
\begin{align*}
\frac{\partial(G^i\circ F)}{\partial x^j}(x)=\sum_{k=1}^m\frac{\partial G^i}{\partial y^k}(F(x))\frac{\partial F^k}{\partial x^j}(x).
\end{align*}
\item 如果 $F$ 与 $G$ 光滑，则 $G\circ F$ 光滑。
\end{enumerate}
\end{corollary}

\begin{exercise}\label{exercise:appC-C-12}
证明推论 C.11。
\end{exercise}

由链式法则与归纳，只要所涉函数充分可微，就可按需要导出复合函数高阶偏导数的公式。

\begin{exercise}\label{exercise:appC-C-13}
假设 $A\subset\mathbb R^n$ 与 $B\subset\mathbb R^m$ 是任意子集，$F:A\to\mathbb R^m$ 与 $G:B\to\mathbb R^p$ 是光滑映射（即在每点的某邻域内有光滑延拓意义下），且 $F(A)\subset B$。证明 $G\circ F:A\to\mathbb R^p$ 光滑。
\end{exercise}

现设 $f:U\to\mathbb R$ 是开子集 $U\subset\mathbb R^n$ 上的光滑实值函数，$a\in U$。对每个向量 $v\in\mathbb R^n$，我们定义 $f$ 在 $a$ 处沿方向 $v$ 的\emph{方向导数}为数
\begin{align}\label{align:appC-directional}
D_v f(a)=\left.\frac{\mathrm{d}}{\mathrm{d}t}\right|_{t=0}f(a+tv).
\end{align}
（此定义对任意向量 $v$ 都有意义；我们不要求 $v$ 是单位向量，正如初等微积分中有时那样。）

由于 $D_v f(a)$ 是复合函数 $t\mapsto a+tv\mapsto f(a+tv)$ 的普通导数，由链式法则它可更具体地写成
\begin{align*}
D_v f(a)=\sum_{i=1}^n v^i\frac{\partial f}{\partial x^i}(a)=\mathrm{D}f(a)v.
\end{align*}

微积分基本定理表达了积分与导数间一个熟知关系。另一个关系是：光滑函数的积分可在积分号下求导。精确陈述见下一定理；这不是所能证明的最佳结果，但对我们的目的已绰绰有余。证明见 [Apo74, Rud76, Str00]。

\begin{theorem}[积分号下求导]\label{theorem:appC-differentiate-under-integral}
设 $U\subset\mathbb R^n$ 是开子集，$a,b\in\mathbb R$，$f:U\times[a,b]\to\mathbb R$ 是连续函数，使对 $i=1,\dots,n$，偏导数 $\partial f/\partial x^i:U\times[a,b]\to\mathbb R$ 在 $U\times[a,b]$ 上存在且连续。用
\begin{align*}
F(x)=\int_a^b f(x,t)\,\mathrm{d}t
\end{align*}
定义 $F:U\to\mathbb R$。则 $F$ 是 $C^1$ 类的，其偏导数可在积分号下求导计算：
\begin{align*}
\frac{\partial F}{\partial x^i}(x)=\int_a^b\frac{\partial f}{\partial x^i}(x,t)\,\mathrm{d}t.
\end{align*}
\end{theorem}

你大概熟悉 Taylor 定理，它表明一个充分光滑的函数如何在一点附近用多项式逼近。我们需要多元 Taylor 定理的一个版本，它给出余项项的显式积分形式。为简洁表达，引入一些简写记号有帮助。对任意指标 $m$ 元组 $I=(i_1,\dots,i_m)$（$1\leqslant i_j\leqslant n$），令 $|I|=m$ 表示 $I$ 中指标个数，并令
\begin{align*}
\partial_I=\frac{\partial^m}{\partial x^{i_1}\cdots\partial x^{i_m}},\qquad
(x-a)_I=(x^{i_1}-a^{i_1})\cdots(x^{i_m}-a^{i_m}).
\end{align*}

\begin{theorem}[Taylor 定理]\label{theorem:appC-taylor}
设 $U\subset\mathbb R^n$ 是开子集，$a\in U$ 固定。假设对某个 $k\geqslant0$ 有 $f\in C^{k+1}(U)$。如果 $W$ 是 $U$ 中含 $a$ 的任意凸子集，则对所有 $x\in W$，
\begin{align}\label{align:appC-taylor}
f(x)=P_k(x)+R_k(x),
\end{align}
其中 $P_k$ 是 $f$ 在 $a$ 处的 $k$ 阶 Taylor 多项式，定义为
\begin{align}\label{align:appC-taylor-poly}
P_k(x)=f(a)+\sum_{m=1}^k\frac{1}{m!}\sum_{I:|I|=m}\partial_I f(a)(x-a)_I,
\end{align}
而 $R_k$ 是 $k$ 阶余项，由
\begin{align}\label{align:appC-taylor-remainder}
R_k(x)=\frac{1}{k!}\sum_{I:|I|=k+1}(x-a)_I\int_0^1(1-t)^k\partial_I f(a+t(x-a))\,\mathrm{d}t
\end{align}
给出。
\end{theorem}

\begin{proof}
对 $k=0$（其中我们把 $P_0$ 理解为 $f(a)$），这只是把微积分基本定理应用于函数 $u(t)=f(a+t(x-a))$，连同链式法则。假设结果对某个 $k$ 成立，对余项中的积分作分部积分，得
\begin{align*}
&\int_0^1(1-t)^k\partial_I f(a+t(x-a))\,\mathrm{d}t\\
&=\left[-\frac{(1-t)^{k+1}}{k+1}\partial_I f(a+t(x-a))\right]_{t=0}^{t=1}
+\int_0^1\frac{(1-t)^{k+1}}{k+1}\frac{\partial}{\partial t}\big(\partial_I f(a+t(x-a))\big)\,\mathrm{d}t\\
&=\frac{1}{k+1}\partial_I f(a)+\frac{1}{k+1}\sum_{j=1}^n(x^j-a^j)\int_0^1(1-t)^{k+1}\frac{\partial}{\partial x^j}\partial_I f(a+t(x-a))\,\mathrm{d}t.
\end{align*}
把此代入\eqref{align:appC-taylor}，便得到把 $k$ 换成 $k+1$ 的类似公式。
\end{proof}

\begin{corollary}\label{corollary:appC-taylor-remainder}
假设 $U\subset\mathbb R^n$ 是开子集，$a\in U$，对某个 $k\geqslant0$ 有 $f\in C^{k+1}(U)$。如果 $W$ 是 $U$ 中含 $a$ 的凸子集，且在 $W$ 上 $f$ 的所有 $k+1$ 阶偏导数绝对值都被常数 $M$ 控制，则对所有 $x\in W$，
\begin{align*}
|f(x)-P_k(x)|\leqslant\frac{n^{k+1}M}{(k+1)!}|x-a|^{k+1},
\end{align*}
其中 $P_k$ 是由\eqref{align:appC-taylor-poly} 定义的 $f$ 在 $a$ 处的 $k$ 阶 Taylor 多项式。
\end{corollary}

\begin{proof}
\eqref{align:appC-taylor-remainder} 右端有 $n^{k+1}$ 项，每项绝对值被 $(1/(k+1)!)|x-a|^{k+1}M$ 控制。
\end{proof}
\end{document}
